% !TEX root = ProjectDescription.tex
\subsection{Thrust 2: \thrusttwo}
\label{sec:thrust2}

\noindent\textbf{Problem Statement.}
% TODO(glitching): 2-4 sentences. Frame defensively (bug-finding), not as attack tooling.
%   - Finding glitching vulnerabilities today is largely manual, target-specific, and slow.
%   - Thrust 1 shows the compiler-introduced surface is large; we need to find the exploitable
%     subset automatically, at scale, before adversaries do.
%   - Challenge: a faithful-yet-tractable fault model, and search that scales to real binaries.

\smallskip
\noindent\textbf{Research Questions.}
\begin{itemize}[leftmargin=*,nosep]
  \item \textbf{RQ2.1:} What fault model(s) faithfully capture real glitching effects (instruction skip, bit flip, data corruption) on compiled embedded code while remaining tractable for automated analysis?
  \item \textbf{RQ2.2:} How can exploitable fault-injection vulnerabilities be discovered \emph{automatically and scalably} in compiled binaries?
  \item \textbf{RQ2.3:} How do we triage and \emph{validate} candidates to separate genuinely exploitable faults from benign ones, ideally confirming on real hardware?
\end{itemize}

\smallskip
\noindent\textbf{Approach Overview.}
% TODO(glitching): one paragraph tying the tasks below to the RQs above.

\subsubsection{\researchtwoone}
Formalizes a tractable fault model capturing real glitching effects (instruction skip, bit flip, data corruption) and defines a binary-level vulnerability condition for automated analysis.
\noindent\textbf{Goal.} % TODO(glitching)
\noindent\textbf{Approach.} % TODO(glitching): formalize fault model(s) and a binary-level vulnerability condition.
\noindent\textbf{Expected Outcome.} % TODO(glitching)

\subsubsection{\researchtwotwo}
Builds an automated search engine combining fault simulation, symbolic execution, and/or fuzzing under the Task~2.1 fault model to discover exploitable sites in compiled binaries at scale.
\noindent\textbf{Goal.} % TODO(glitching)
\noindent\textbf{Approach.} % TODO(glitching): automated search -- fault simulation / symbolic execution / fuzzing under the fault model.
\noindent\textbf{Expected Outcome.} % TODO(glitching)

\subsubsection{\researchtwothree}
Ranks candidate vulnerabilities by exploitability and confirms high-priority findings on a physical hardware glitching testbed, establishing ground truth for Thrusts~1 and~3.
\noindent\textbf{Goal.} % TODO(glitching)
\noindent\textbf{Approach.} % TODO(glitching): exploitability triage; validation on a hardware glitching testbed.
\noindent\textbf{Expected Outcome.} % TODO(glitching)
